% LaTeX-ready experimental setup. Add the relevant package imports in the manuscript preamble.
\section{Experimental Setup}
\label{sec:experimental-setup}

The benchmark considers static global path planning in known two-dimensional
workspaces. All maps have continuous Cartesian bounds of $100\times100$ units
and contain circular and axis-aligned rectangular obstacles. The robot is
modelled as a disc and collision checking accounts for the specified robot
radius. Nine planners were evaluated: A*, Dijkstra, RRT, RRT*, Informed RRT*,
PRM, ACO, ACO--GA, and the proposed Three-Phase Incremental SFO. Each planner
retained its native representation and internal objective function; common
metrics were calculated only after path generation.

Five independent runs were executed for every map--planner pair using a
deterministic seed schedule initialized from base seed 2026. Hence, the study
contains $9\times3\times5=135$ planned runs.

\begin{table*}[t]
\centering
\caption{Static-map specifications used in the benchmark.}
\label{tab:map-specifications}
\begin{tabular}{lcccccl}
\toprule
Map & Bounds & Start & Goal & Radius & Obstacles & Purpose \\
\midrule
Dense-obstacle & $100\times100$ & $(5,5)$ & $(95,95)$ & 0.5 & $3$ circles, $2$ rectangles & Mixed clutter \\
Narrow-passage & $100\times100$ & $(8,12)$ & $(92,88)$ & 1.0 & $4$ circles, $6$ rectangles & Central passage \\
Bug-trap & $100\times100$ & $(48,50)$ & $(92,50)$ & 1.0 & $3$ circles, $5$ rectangles & U-shaped escape \\
\bottomrule
\end{tabular}
\end{table*}

\begin{table*}[t]
\centering
\caption{Planner settings used in the benchmark.}
\label{tab:planner-settings}
\small
\begin{tabular}{lp{11.2cm}}
\toprule
Planner & Settings \\
\midrule
A* & Grid resolution $=1.0$. \\
Dijkstra & Grid resolution $=1.0$. \\
RRT & Rectangle resolution $=1.0$; expansion distance $=4.0$; path resolution $=0.5$; goal sampling $=10\%$; maximum iterations $=1000$. \\
RRT* & RRT settings with connection-circle distance $=50.0$ and search until the maximum iteration enabled. \\
Informed RRT* & Rectangle resolution $=1.0$; expansion distance $=2.0$; goal sampling $=10\%$; maximum iterations $=1000$. \\
PRM & Rectangle resolution $=1.0$; sampled nodes $=500$; nearest neighbours $=10$; maximum edge length $=30.0$. \\
ACO & Ants $=20$; steps per ant $=200$; colony iterations $=40$; $\alpha=1.1$; $\beta=12.0$; $\gamma=0.02$; evaporation $=0.5$; brushfire iterations $=5$; $Q=10.0$. \\
ACO--GA & Grid resolution $=1.0$; population $=80$; ACO iterations $=90$; GA iterations $=60$; $\rho=0.25$; $\alpha=2.0$; $\beta=8.0$; $Q=100.0$; crossover $=0.20$; mutation $=0.05$; length and smoothness weights $=0.50$; corner cutting disabled; timeout $=1800$ s. \\
Three-Phase Incremental SFO & Population $=30$; iterations $=250$; candidate attempts $=100$; initial step range $=[8.0,16.0]$; glide velocity range $=[5.0,18.0]$; target velocity range $=[0.5,8.0]$; micro-adjustment velocity range $=[0.1,2.5]$; goal-connection distance $=18.0$; collision resolution $=0.25$. \\
\bottomrule
\end{tabular}
\end{table*}

A run was classified as successful only if it reached the goal, was
collision-free, and passed the common geometric validator. Shortest-path,
average-path, turning-point, and smoothness summaries were calculated from
valid completed paths only. Execution time was averaged over all attempts,
including unsuccessful attempts. Failed and partial paths were retained for
diagnostic figures and success-rate analysis but were excluded from path-quality
means.
